\par\textbf{Procedimiento completo para comparar tu solución.} El combustible añade caudal a la salida. Se conserva esa masa tanto en el balance de momento como en la energía cinética. La eficiencia global compara potencia útil con aporte químico; no se identifica con rendimiento propulsivo.

\begin{enumerate}

\item Caudal de combustible por relación combustible/aire.

\[\fitmath{\mathrm{mf}=0.02\cdot 25\approx 0.5\,\mathrm{kg/s}}\]

\item Caudal total saliente.

\[\fitmath{\mathrm{mg}=25+\mathrm{mf}}\]

\[\fitmath{\mathrm{mg}=25+(0.5)\approx 25.5\,\mathrm{kg/s}}\]

\item Balance axial de momento.

\[\fitmath{F=\mathrm{mg}\cdot 600-25\cdot 200}\]

\[\fitmath{F=(25.5)\cdot 600-25\cdot 200\approx 10300\,\mathrm{N}}\]

\item Potencia útil al vuelo.

\[\fitmath{\mathrm{Pu}=F\cdot 200}\]

\[\fitmath{\mathrm{Pu}=(10300)\cdot 200\approx 2060000\,\mathrm{W}}\]

\item Aumento de energía cinética de las corrientes.

\[\fitmath{\mathrm{Ek}=\frac{\mathrm{mg}\cdot {600}^{2}-25\cdot {200}^{2}}{2}}\]

\[\fitmath{\mathrm{Ek}=\frac{(25.5)\cdot {600}^{2}-25\cdot {200}^{2}}{2}\approx 4090000\,\mathrm{W}}\]

\item Aporte químico.

\[\fitmath{\mathrm{Pchem}=\mathrm{mf}\cdot 43000000}\]

\[\fitmath{\mathrm{Pchem}=(0.5)\cdot 43000000\approx 2.15\times10^{7}\,\mathrm{W}}\]

\item Eficiencia global.

\[\fitmath{\mathrm{etao}=\frac{\mathrm{Pu}}{\mathrm{Pchem}}}\]

\[\fitmath{\mathrm{etao}=\frac{(2060000)}{(21500000)}\approx 0.09581395\,}\]

\item Consumo por empuje.

\[\fitmath{\mathrm{TSFC}=\frac{\mathrm{mf}}{F}}\]

\[\fitmath{\mathrm{TSFC}=\frac{(0.5)}{(10300)}\approx 4.854369\times10^{-5}\,\mathrm{kg/(N\,s)}}\]

\item Comparación que desprecia masa de combustible.

\[\fitmath{\mathrm{Fapprox}=25\cdot \left(600-200\right)\approx 10000\,\mathrm{N}}\]

\end{enumerate}
